\section{Introduction}
\label{sec:introduction}

Benchmark leaderboards report accuracy as if it were observed. It is not: a
model's reported score is the output of a fixed, deterministic verifier
applied to free-text generations, and that verifier can be wrong. When it
is, the leaderboard number is the model's accuracy contaminated by the
verifier's own error rate, in a direction and magnitude nobody has
measured (schematic: Appendix~\ref{sec:appendix-figures},
Figure~\ref{fig:conceptual}).

This paper asks what a benchmark comparison is still entitled to claim once
that contamination is taken seriously. We audit MATH-Hard's standard
comparator on real model responses, find it disagrees with human judgment
at a small but nonzero rate, and derive the set of true accuracy values
consistent with what was observed---not a point estimate corrected by a
guessed error rate, but a \emph{partially identified set}, sharp under the
weakest assumption the audit data will bear. The central finding: audit
uncertainty, not sample size, is usually the binding constraint. Audit-
estimation uncertainty---how few human-labeled items exist to pin down the
verifier's error rate---exceeds the field's already-reported sampling
uncertainty for the majority of real model comparisons we can identify at
all.

That finding raises a design question the literature has not answered: if
verifier auditing is the binding constraint, how should a fixed
human-labeling budget be spent? We formalize this as a stratified-sampling
allocation problem, derive the classical \citep{neyman1934stratified}
solution, and show---on our own real, sparse audit data---that the
classical formula's assumption (a decent pilot estimate) fails exactly
where most benchmarks are actually audited: when most strata have observed
zero error events, naive plug-in allocation concentrates budget on
whichever cell drew a noisy nonzero rate, and is dominated by simply
spreading the budget evenly. A stress-tested, sparse-data-aware allocation
rule is this paper's most directly usable output.

A third thread: MATH-Hard's evaluation harness computes accuracy two
different ways from the same generations---LaTeX-aware symbolic parsing
(\exactmatch) and last-\code{\textbackslash boxed\{\}} matching
(\scoreboxed)---and these disagree on 3.63\% of real responses. Every human
audit here was conducted against \scoreboxed; the field's reported accuracy
comes from \exactmatch. This mismatch does not overturn the
audit-estimation-dominance finding, tested directly including under an
adversarial worst-case bound, but establishes a second, independent source
of benchmark uncertainty: \emph{which} automated criterion scores a
benchmark is itself a measurable, consequential choice, independently
corroborated by production-scale evidence (Section~\ref{sec:related-work}).

\paragraph{Contributions.} (1) A sharp partial-identification treatment of
benchmark accuracy under an audited, imperfect verifier, adapting classical
misclassification-bounds machinery
\citep{manski1995identification,molinari2008partial} to benchmark
comparison. (2) An empirical demonstration that audit-estimation
uncertainty dominates sampling uncertainty for the majority of identifiable
comparisons on a widely-used benchmark. (3) A stratified audit-design
methodology (Section~\ref{sec:audit-design}), including a documented
failure mode of the textbook-optimal allocation rule under sparse real data
and a tested alternative. (4) A secondary, honestly-scoped finding that
automated-scorer choice is itself a measurable, non-dominant source of
benchmark uncertainty, with a worst-case bound rather than an assumed-safe
headline.
